\documentclass[12pt]{article}
\usepackage{amsfonts,amsthm,amsmath,amssymb,mathtools}
\usepackage{mathrsfs}
\usepackage{enumitem}
\usepackage{tikz-cd}
\usepackage{geometry}
\usepackage{mlmodern}
\usepackage{xcolor}

\geometry{letterpaper, portrait, margin=1in}

\setcounter{section}{0}

\renewcommand{\thesection}{\S\arabic{section}}
\renewcommand{\thesubsection}{\arabic{section}}

\DeclareMathOperator{\lcm}{lcm}
\DeclareMathOperator{\im}{im}

\newtheorem{thm}{Theorem}
\newenvironment{theorem}[1]{
	\renewcommand{\thethm}{#1}
	\begin{thm}
		}{\end{thm}}

\newtheorem{lma}{Lemma}
\newenvironment{lemma}[1]{
	\renewcommand{\thelma}{#1}
	\begin{lma}
		}{\end{lma}}

\theoremstyle{definition}
\newtheorem{ex}{}
\newenvironment{exercise}[1]{
	\renewcommand{\theex}{#1}
	\begin{ex}
		}{\end{ex}}

\title{Homework 3}
\author{Connor Mooneyhan}
\date{September 22, 2025}

\begin{document}

\maketitle

\begin{ex}
	Let $G$ be a group and let $H$ be a subgroup of $G$.
	Prove that $C_G(H) \subset N_G(H)$.
	Prove that equality holds if $H$ has order two.
\end{ex}
\begin{proof}
	If $g \in C_G(H)$, then $ghg^{-1} = h$ for each $h \in H$, so $gHg^{-1} = H$.
	Thus $g \in N_G(H)$.

	If $H$ has order two, then one of the two elements of $H$ is $e$, and the other is an element of order two, say $x$.
	If $g \in N_G(H)$, then $geg^{-1} = e$, and $g$ must act on $x$ to take it to the only remaining element of $H$, which is $x$ itself, so $gxg^{-1} = x$.
\end{proof}

\begin{ex}
	Let $G$ be a group, and suppose $H$ is a finite subgroup of $G$ generated by a set $X$.
	Prove that $g \in N_G(H)$ if and only if $gxg^{-1} \in H$ for all $x \in X$.
	Conclude that for $x \in G$, one has that $g \in N_G(\left\langle x \right\rangle)$ if and only if $gxg^{-1} = x^a$ for some $a \in \mathbb{Z}$.
\end{ex}
\begin{proof}
	If $g \in N_G(H)$, then $gHg^{-1}$, so indeed $gxg^{-1} \in H$ for any $x \in X \subset H$.

	Now let $g \in N_G(H)$ and assume that $gxg^{-1} \in H$ for all $x \in X$.
	Then given $h \in H$, we can write $h = gh' g^{-1}$ for some $h' \in H$, so $h' = x_1^{a_1} \cdots x_n^{a_n}$ for some $x_i \in X$ and each $a_j \in \mathbb{Z}$.
	We then write \begin{equation*}
		h = gh'g^{-1} = gx_1^{a_1}x_2{a_2}\cdots x_n^{a_n}g^{-1} = (gx_1g^{-1})^{a_1}(gx_2g^{-1})^{a_2}\cdots (gx_ng^{-1})^{a_n},
	\end{equation*}
	so $h$ is a product of elements of $H$.
	Thus $gHg^{-1} \subset H$.

	Now assume that $h, h' \in gHg^{-1}$.
	Then \begin{equation*}
		ghg^{-1} = gh'g^{-1} \implies h = h'.
	\end{equation*}
	Thus since $H$ is finite and $f: H \to H$ defined by $f(h) = ghg^{-1}$ is injective, $f$ is in fact a bijection, and $gHg^{-1} = H$.
\end{proof}

\begin{ex}
	Let $H = \left\langle \begin{pmatrix}1 & 1 \\ 0 & 1\end{pmatrix} \right\rangle \subset \mathrm{GL}_2(\mathbb{R})$, and let $g = \begin{pmatrix}2 & 0 \\ 0 & 1\end{pmatrix}$.
	Show that $gHg^{-1} \subset H$, but that this containment is proper and hence $g$ is \textit{not} in the normalizer of $H$.
	Why does this not violate the statement in Problem 2?
\end{ex}
\begin{proof}
	Since any element of $H$ is of the form $\begin{pmatrix}1 & 1 \\ 0 & 1\end{pmatrix}^n = \begin{pmatrix}1 & n \\ 0 & 1\end{pmatrix}$ and $g^{-1} = \begin{pmatrix}\frac{1}{2} & 0 \\ 0 & 1\end{pmatrix}$, we can calculate explicitly \begin{equation*}
		\begin{pmatrix}2 & 0 \\ 0 & 1\end{pmatrix} \begin{pmatrix} 1 & n \\ 0 & 1\end{pmatrix} \begin{pmatrix}\frac{1}{2} & 0 \\ 0 & 1\end{pmatrix} = \begin{pmatrix}2 & 2n \\ 0 & 1\end{pmatrix}\begin{pmatrix}\frac{1}{2} & 0 \\ 0 & 1\end{pmatrix} = \begin{pmatrix}1 & 2n \\ 0 & 1\end{pmatrix}.
	\end{equation*}
	Since this only includes matrices with an even upper-right entry, it excludes $\begin{pmatrix}1 & 1 \\ 0 & 1\end{pmatrix}$ itself, and thus $gHg^{-1} \subsetneq H$.
	This is in complete agreement with Problem 2, since Problem 2 required the subgroup in question, here $H$, to be finite, and $H$ is infinite.
\end{proof}

\begin{ex}
	Let $G$ be a group, and $n$ a positive integer.
	Prove that the set of group homomorphisms from $\mathbb{Z}_n$ to $G$ are in bijection with the set of elements of $G$ of order dividing $n$.
\end{ex}
\begin{proof}
	We will write $\mathbb{Z}_n = \left\langle x : x^n \right\rangle$ for ease of notation.
	The idea here is that a homomorphism from $\mathbb{Z}_n$ to $G$ is uniquely determined by where it sends $x$.
	Indeed, given a homomorphism $\varphi$ which sends $x$ to $g \in G$, then $\varphi(x^i) = \varphi(x)^i = g^i$, so $\varphi$ is described completely by where it sends $x$ since $x$ generates $\mathbb{Z}_n$.
	A direct consequence is the injectivity we need.
	Then, using the fact that $|\varphi(x)|$ divides $|x|=n$, we determine that we can only send $x$ to elements of $G$ with order dividing that of $x$.
	In fact, given any $g \in G$ with $|g|$ dividing $n$, there is a homomorphism determined entirely by sending $x$ to $g$, so we have surjectivity.
\end{proof}

\begin{ex}
	A permutation $\sigma \in S_n$ of the form $(i\;j)$ for $1 \leq i < j \leq n$ (in cycle notation) is called a \textit{transformation}.
	If $j = i + 1$, we say the transposition is \textit{adjacent}.
	Prove that $S_n$ is generated by the set of transpositions.
	Then prove that all transpositions are products of adjacent transpositions to conclude that $S_n$ is in fact generated by adjacent transpositions.
\end{ex}
\begin{proof}
	First, we show that every element of $S_n$ is a product of transpositions.
	Indeed, for a cycle $(a_1\;a_2\;\cdots\;a_m)$, we can write \begin{equation*}
		(a_1\;a_2\;\cdots\;a_m) = (a_{m-1}\;a_m)\cdots(a_2\;a_m)(a_1\;a_m).
	\end{equation*}

	Given any transposition, then, we want to show that it is a product of adjacent transpositions, which we will prove by induction.
	We write any transposition as $(a\;[a+i])$.
	Let $\sigma \in S_n$.
	If $i = 0$, then $\sigma$ is the identity permutation, which is the product of zero adjacent transpositions.
	If $i = 1$, Then $\sigma = (a\;[a+1])$ is an adjacent permutation by definition.
	Now assume that $i=n$ and the hypothesis holds for $i < n$.
	Then we can write \begin{align*}
		(a\;[a+n]) = (a\;[a+n-1])([a+n-1]\;[a+n])(a\;[a+n-1]),
	\end{align*}
	and since $(a\;[a + n - 1])$ can be written as a product of adjacent permutations, this represents $(a\;[a+n])$ as a product of adjacent permutations.
\end{proof}

\begin{ex}
	Let $p \geq 3$ be a prime integer.
	Prove that $D_{2p}$ is generated by any pair of distinct reflections.
	Show by means of example that this can fail if $p$ is not prime.
\end{ex}
\begin{proof}
	Choose a reflection, which we call $s$, and a rotation, which we call $r$.
	We can now choose a reflection distinct from $s$ that can be written as $r^ns$ for some $n$.

	First, we observe that since $p$ is prime, the subgroup generated by $r$ has order $p$, regardless which rotation we have chosen.
	We now want to show that $s$ and $r^ns$ generate the group.
	Note that $s$ has order 2, and since $(r^ns)(s) = r^n$ is an element of the subgroup generated by $r$, we have that $r^{nm} = r$ for some $m$.
	Thus we have produced both standard generators of $D_{2n}$ ($r$ and $s$), and all relations are satisfied ($r^p = s^2 = e$ together with the inherited $r^is = sr^{-i}$ for any $i$).
\end{proof}

\begin{ex}
	Let $G$ be a group and let $H$ be a subgroup.
	Let $G$ act on the set $G/H$ by left translation.
	Prove that the kernel of this action is $K = \bigcap_{g \in G} gHg^{-1}$, and prove that $K$ is the largest normal subgroup of $G$ contained in $H$.
\end{ex}
\begin{proof}
	Let $k \in K$.
	First we show that $K \subset \bigcap_{g \in G} gHg^{-1}$.
	We want to prove that given $g \in G$, $k = ghg^{-1}$ for some $h \in H$, so we will show that $g^{-1}kg \in H$, since that would suffice for $h$.
	Since $k \in K$, $kgH = gH$, so $g^{-1}kgH = g^{-1}(kgH) = (g^{-1}g)H = H$.
	Thus, $g^{-1}kg \in H$ as desired.
	For the reverse containment, fix $g \in G$ and let $k = ghg^{-1}$ for some $h \in H$.
	Then $ghg^{-1}(gH) = gh(g^{-1}g)H = ghH = gH$, so since $g$ was arbitrary, $ghg^{-1} \in K$.

	Now let $J$ be a normal subgroup of $G$ contained in $H$, and let $j \in J$.
	Then given any $g \in G$, since $gJg^{-1} = J$, there exists a $j' \in J$ such that $gj'g^{-1} = j$.
	Since $j' \in J \subset H$ and the choice of $g$ was arbitrary, we have \begin{equation*}
		j \in \bigcap_{g \in G}gHg^{-1} = K \implies J \subset K.
	\end{equation*}
\end{proof}

\begin{ex}
	Suppose $G$ acts on a set $X$, let $x \in X$, and let $H = \mathrm{Stab}_G(x)$.
	Recall that $G/H$ denotes the set of left cosets of $H$ in $G$.
	Prove that the assignment \begin{align*}
		\varphi: G/H & \to \mathrm{Orb}_G(x) \\
		gH           & \mapsto g\cdot x
	\end{align*}
	is a well-defined bijective function.
	Furthermore, show that $\varphi$ is $G$-\textit{equivalent}; that is, prove that for all $a \in G$ and all $gH \in G/H$, one has that $\varphi(a \cdot gH) = a \cdot \varphi(gH)$, where the $G$-action on $G/H$ is left translation of cosets, and the $G$-action on $\mathrm{Orb}_G(x)$ is inherited from the action of $G$ on $X$.
\end{ex}
\begin{proof}
	Let $g_1, g_2 \in G$ such that $g_1H = g_2H$.
	Then given $h_1 \in H$, there exists an $h_2 \in H$ such that $g_1h_1 = g_2h_2$.
	We can then deduce that \begin{align*}
		g_1 \cdot x & = g_1 \cdot (h_1 \cdot x) \\
		            & = (g_1h_1) \cdot x        \\
		            & = (g_2h_2) \cdot x        \\
		            & = g_2 \cdot (h_2 \cdot x) \\
		            & = g_2 \cdot x,
	\end{align*}
	so $\varphi$ is a well-defined function.

	By a similar line of argument, we can show injectivity.
	Assume $g_1H \neq g_2H$ for some $g_1, g_2 \in G$.
	Then $g_1H$ and $g_2H$ are disjoint, so given any $h_1, h_2 \in H$, $g_1h_1 \neq g_2h_2$.
	Then we have \begin{align*}
		g_1 \cdot x & = g_1 \cdot (h_1 \cdot x) \\
		            & = (g_1h_1) \cdot x        \\
		            & \neq (g_2h_2) \cdot x     \\
		            & = g_2 \cdot (h_2 \cdot x) \\
		            & = g_2 \cdot x.
	\end{align*}
	Finally, given any $g \cdot x \in \mathrm{Orb}_G(x)$, the coset $gH$ maps to $g\cdot x$, so $\varphi$ is a bijective function.

	We conclude by showing $G$-equivariance: \begin{align*}
		\varphi(a \cdot gH) & = \varphi((ag)H) \\
                        & = (ag) \cdot x \\
                        & = a \cdot (g \cdot x) \\
                        & = a \cdot \varphi(gH).
	\end{align*}
\end{proof}

\begin{ex}
	Let $\varphi: G \to H$ be a group homomorphism.
	Prove that if $N$ is a subgroup of $H$ then $\varphi^{-1}(N)$ is a subgroup of $G$.
	Prove that if $N$ is normal, so is $\varphi^{-1}(N)$.
\end{ex}
\begin{proof}
	Let $g, h \in \varphi^{-1}(N)$.
	Then $\varphi(gh^{-1}) = \varphi(g)\varphi(h)^{-1} \in N$, so $gh^{-1} \in \varphi^{-1}(N)$.

	Now given $g \in G$ and $n \in \varphi^{-1}(N)$, we have $\varphi(gng^{-1}) = \varphi(g)\varphi(n)\varphi(g)^{-1} \in N$ since $\varphi(n) \in N$ and $N$ is normal.
	Thus $gng^{-1} \in \varphi^{-1}(N)$, so since $g$ and $n$ were arbitrary elements of $G$ and $\varphi^{-1}(N)$ respectively, $\varphi^{-1}(N)$ is normal.
\end{proof}

\begin{ex}
	Suppose that $G$ is a group and that $X$ is a subset of $G$ that is closed under conjugation.
	Prove that $\left\langle X \right\rangle$ is a normal subgroup.
	Next, let $Y$ be the set of commutators \begin{equation*}
		\left\lbrace aba^{-1}b^{-1} : a, b \in G \right\rbrace
	\end{equation*}
	and prove that $Y$ is closed under conjugation.
	Use these facts to show that the commutator subgroup $[G,G] = \left\langle Y \right\rangle$ is a normal subgroup.
\end{ex}
\begin{proof}
	If $n \in \left\langle X \right\rangle$ and $g \in G$, write $n=n_1n_2 \cdots n_m$ where each $n_i \in X$.
	Then \begin{equation*}
		gng^{-1} = gn_1\cdots n_m g^{-1} = (gn_1g^{-1})\cdots(gn_mg^{-1}) \in \left\langle X \right\rangle
	\end{equation*}
	since $gn_ig^{-1} \in X$ for each $i$.

	Now, given any $a, b, g \in G$, we have \begin{equation*}
		g(aba^{-1}b^{-1})g^{-1} = (gag^{-1})(gbg^{-1})(gag^{-1})^{-1}(gbg^{-1})^{-1} \in Y.
	\end{equation*}
	Thus $[G,G] = \left\langle Y \right\rangle$ is a normal subgroup.
\end{proof}

\end{document}
